\begin{definition}[The Gaussian Sobolev space $W^{1,2}(\gamma_n)$]
Set
\[
  W^{1,2}(\gamma_n)
    = \bigl\{ v \in L^2(\gamma_n) :
        \text{$v$ has a weak Euclidean gradient $\nabla v$ with }
        2^{-1/2}\nabla v \in L^2(\gamma_n;\mathbb{R}^n) \bigr\},
\]
the weak gradient being understood in the sense of \noderef{HasGaussWeakGrad} and
$L^2(\gamma_n;\mathbb{R}^n)$ denoting the space of $\mathbb{R}^n$-valued functions
square integrable against the measure $\gamma_n$ of \noderef{GaussMeasure}, with
the Euclidean norm on the target.  Equivalently, $v \in W^{1,2}(\gamma_n)$ iff
$v \in L^2(\gamma_n)$ and there is $\Gamma \in L^2(\gamma_n;\mathbb{R}^n)$ with
$\sqrt2\,\Gamma$ a weak Euclidean gradient of $v$; the factor $\sqrt2$ records the
energy normalization $D = 2^{-1/2}\nabla$ of the reference paper, so that $\Gamma$
is $Dv$.  This is the characterization of $W^{1,2}(\gamma_n)$ used throughout: the
functions in $L^2(\gamma_n)$ whose weak Euclidean derivatives again lie in
$L^2(\gamma_n)$.  The reference paper notes that this coincides with the
completion of $C_c^\infty(\mathbb{R}^n)$ in the norm
$\|v\|^2 = \|v\|_2^2 + \tfrac12\|\nabla v\|_2^2$; that equivalence is not used
anywhere below, and only the weak-derivative description is taken as the
definition.
\end{definition}
